\documentclass[11pt,a4paper]{article}
\usepackage[utf8]{inputenc}
\usepackage[T1]{fontenc}
\usepackage{lmodern}
\usepackage[margin=1in]{geometry}
\usepackage{hyperref}
\usepackage{longtable}
\usepackage{array}
\usepackage{enumitem}
\usepackage{amsmath,amssymb}

\title{Fractional Semantics Project\\Implementation Report}
\author{Project Workspace: Fractional Semantics}
\date{\today}

\begin{document}
\maketitle
\tableofcontents
\bigskip

\section{Goal}
The project was developed to compute Fractional Semantics values from sequents, handle a belief set $B$, support belief revision workflows, and generate readable decomposition outputs in classic prooftree style (with Fractional decorations).

The final system includes:
\begin{itemize}[leftmargin=1.5em]
  \item CLI workflows (analysis, revision, demos, certificates, comparison tools),
  \item automatic \texttt{.tex}/\texttt{.pdf} output folders for each run,
  \item a browser-based visual interface with decomposition view, PNG preview, history, and export.
\end{itemize}

\section{Core Architecture}
The implementation is modular. Main components:

\begin{longtable}{>{\raggedright\arraybackslash}p{0.34\linewidth} >{\raggedright\arraybackslash}p{0.60\linewidth}}
\textbf{Module} & \textbf{Responsibility} \\
\hline
\texttt{fs\_semantics/logic.py} & Formula AST, parser, one-sided/two-sided sequent parsing, rendering with precedence-aware parentheses. \\
\texttt{fs\_semantics/calculus.py} & GS4-like decomposition and Fractional scoring on leaves. \\
\texttt{fs\_semantics/beliefs.py} & Belief base parsing and cut-closure support. \\
\texttt{fs\_semantics/latex.py} & Classic \texttt{bussproofs} decomposition rendering with \texttt{\textbackslash sststile\{\textbackslash mathrm\{n\}\}\{m\}}. \\
\texttt{fs\_semantics/latex\_compile.py} & LaTeX to PDF/PNG compilation pipeline, landscape tree-only previews, cropping margins. \\
\texttt{fs\_semantics/decomposition\_service.py} & Orchestrates parsing, decomposition, tree payload, validation, and output bundle. \\
\texttt{fs\_semantics/visualization.py} & JSON payload for visual tree panel (rule labels, $N/M$, display sequent). \\
\texttt{fs\_semantics/proof\_validation.py} & Structural validation of decomposition steps and root coherence. \\
\texttt{fs\_semantics/propositional\_sc.py} & Pure propositional sequent-calculus reference engine for decomposition comparison. \\
\texttt{fs\_semantics/afp\_certificate.py} & Isabelle/AFP-oriented certificate generation from produced proof trees. \\
\texttt{fs\_semantics/web\_gui.py} & Browser UI (form, tree, decomposition, PNG preview, history, random sequent generation). \\
\texttt{fs\_semantics/cli.py} & Command-line entrypoints and end-to-end workflows. \\
\end{longtable}

\section{Input Semantics and Normalization}
\subsection{Supported Sequent Forms}
Two forms are accepted:
\begin{itemize}[leftmargin=1.5em]
  \item One-sided: \texttt{A, B, C}
  \item Two-sided: \(\Gamma \vdash \Delta\), \texttt{Gamma |- Delta}, or \texttt{Gamma \textbackslash vdash Delta}
\end{itemize}

\subsection{Comma Meaning}
\begin{itemize}[leftmargin=1.5em]
  \item Left side commas are conjunction.
  \item Right side commas are disjunction.
  \item Two-sided sequents are normalized to one-sided form:
  \[
    \Gamma \vdash \Delta \quad \mapsto \quad \vdash \neg\Gamma,\Delta
  \]
\end{itemize}

\subsection{Parentheses Discipline}
Rendering is precedence-aware and keeps only necessary parentheses. The system was refined to avoid redundant bracket pairs (for example avoiding artificial forms like double wrapping from rendering artifacts), while preserving logical correctness.

\section{Fractional Decomposition and Belief Modes}
\subsection{Classic Prooftree Output}
The decomposition is generated in \texttt{bussproofs} style (not analytical textual decomposition), decorated with Fractional turnstiles:
\[
\texttt{\textbackslash sststile\{\textbackslash mathrm\{N\}\}\{M\}}
\]
with rule labels such as \texttt{($\vee$)}, \texttt{($\wedge$)}, \texttt{(ax.)}, \texttt{(b)}, \texttt{($\overline{\text{ax.}}$)}.

\subsection{Belief Set Handling}
Beliefs can be provided by repeated options or prompt-style text:
\begin{itemize}[leftmargin=1.5em]
  \item \texttt{--belief "(p | q)" --belief "\textasciitilde u"}
  \item \texttt{--belief-set "B=\{(p | q); \textasciitilde u\}"}
\end{itemize}

If $B$ is empty, the decomposition uses undecorated turnstiles (no $\,\_\{\mathbb{B}\}$ suffix). If $B$ is non-empty, belief-sensitive values are enabled.

\subsection{Full Beliefs Mode}
A dedicated full-beliefs mode was added to include explicit belief contributions:
\[
\delta_{\{\cdots\}}
\]
inside the $M$ component, so each belief contribution is visible in the decomposition.

\section{Belief Revision Workflows}
The CLI supports:
\begin{itemize}[leftmargin=1.5em]
  \item contraction,
  \item expansion,
  \item revision (Levi identity),
  \item automatic revision demo generation.
\end{itemize}

Each run produces a new timestamped folder in \texttt{tex\_outputs/} and compiles generated LaTeX automatically.

\section{Output Pipeline and LaTeX Compilation}
\subsection{Run Isolation}
Every new computation creates a dedicated subfolder under \texttt{tex\_outputs/} so files do not overwrite each other.

\subsection{Preview Compilation}
For GUI preview:
\begin{itemize}[leftmargin=1.5em]
  \item decomposition is compiled in a landscape-oriented LaTeX preview document,
  \item tree-only PNG output is generated (PDF-to-PNG conversion toolchain),
  \item DPI is configurable for high-resolution output,
  \item border padding is applied to keep extreme formulas readable.
\end{itemize}

\subsection{Custom Scaled Prooftree Environment}
The following environment is integrated in preview generation to control render scale:
\begin{verbatim}
\newenvironment{scprooftree}[1]%
{\gdef\scalefactor{#1}\begin{center}\proofSkipAmount \leavevmode}%
{\scalebox{\scalefactor}{\DisplayProof}\proofSkipAmount \end{center} }
\end{verbatim}

\section{Visual Interface (Web GUI)}
\subsection{Why Web GUI}
When \texttt{tkinter} is unavailable in the local Python build, the system uses a browser-based interface as fallback (and can also be launched directly).

\subsection{Main Features Implemented}
\begin{itemize}[leftmargin=1.5em]
  \item Inputs: sequent, belief set, belief mode, DPI.
  \item Outputs: decomposition-only LaTeX and interactive visual tree.
  \item PNG preview area with scroll and zoom controls.
  \item Loading indicator during decomposition and compilation.
  \item Legend for OR/AND/negation accepted syntax.
  \item History dropdown with restoration of previous runs.
  \item PNG download button.
  \item Random sequent button.
\end{itemize}

\subsection{Random Sequent Generation}
A random AST-based formula generator was added so produced sequents are parser-safe and syntactically valid. Both one-sided and two-sided forms are generated.

\subsection{Preview Usability Fixes}
Several UI fixes were introduced:
\begin{itemize}[leftmargin=1.5em]
  \item preview recenters automatically when image loads,
  \item default preview readability improved,
  \item zoom strategy adjusted for large trees.
\end{itemize}

Final zoom logic is based on decomposition complexity by counting \texttt{\textbackslash BinaryInfC} occurrences, e.g.:
\begin{itemize}[leftmargin=1.5em]
  \item 0 binary rules $\rightarrow$ 120\%
  \item 1 binary rule $\rightarrow$ 110\%
  \item 3 binary rules $\rightarrow$ 55\%
\end{itemize}
with clamped behavior for larger trees.

\section{Correctness Controls}
\subsection{Structural Validation}
Each generated proof tree is checked for:
\begin{itemize}[leftmargin=1.5em]
  \item correct parent/child rule behavior for $\vee$ and $\wedge$ nodes,
  \item expected normalized root sequent,
  \item coherent decomposition under formula permutations where appropriate.
\end{itemize}

\subsection{Parentheses-Aware Consistency}
Validation and rendering were refined to keep output aligned with input grouping semantics, including correct handling of right-side comma as disjunction and left-side comma as conjunction.

\section{Reference Comparison and AFP Alignment}
\subsection{Pure Sequent-Calculus Comparator}
A dedicated propositional sequent-calculus reference module was added to compare decomposition shapes against the Fractional implementation (while keeping Fractional decorations separate).

\subsection{ATP/AFP Coherence}
The project remains intentionally restricted to Fractional Semantics and its GS4-style fragment, but includes:
\begin{itemize}[leftmargin=1.5em]
  \item ATP-coherence checks in decomposition logic,
  \item optional AFP/Isabelle certificate generation,
  \item proof-shape checking against imported sequent-calculus infrastructure.
\end{itemize}

\section{Testing}
Regression tests were expanded over time to cover:
\begin{itemize}[leftmargin=1.5em]
  \item parser semantics and comma interpretation,
  \item decomposition root correctness,
  \item LaTeX/PDF/PNG compilation constraints,
  \item proof validation logic,
  \item SC comparison behavior,
  \item AFP certificate generation,
  \item random sequent generation validity.
\end{itemize}

At report generation time, the test suite in this workspace runs successfully.

\section{Result}
The delivered system now provides:
\begin{itemize}[leftmargin=1.5em]
  \item Fractional Semantics computation from sequents,
  \item belief-aware decomposition in classic prooftree format,
  \item standard and full-beliefs modes (with $\delta$ decorations),
  \item automatic run-folder generation and compilation,
  \item an operational visual interface for decomposition exploration,
  \item validation and comparison tooling to support proof reliability.
\end{itemize}

\end{document}
